% !TeX program = lualatex
% halo:begin
% {
%   "version": 1,
%   "publish": true,
%   "course": "higher-algebra-i",
%   "section": "2.1",
%   "title": "高等代数（一）2.1：线性相关性与向量组的秩",
%   "categories": ["study-notes"],
%   "tags": ["mathematics"],
%   "excerpt": "线性相关与线性无关、等价向量组、向量个数比较、极大线性无关组及向量组的秩。"
% }
% halo:end
% 课堂来源：IMG_2689--2697.jpeg；照片未提供可确认的课堂日期。
% 续整理：IMG_3181左页，北大《高等代数》第六版习题III第13题，印刷页107/PDF第112页。
% 2.1为笔记自身的主题编号，不是教材节号或课堂讲次。
% 教材：北大《高等代数》第五版第三章，n维向量空间与线性相关性。
% 所用PDF：/Users/yuchenzhang/Downloads/高等代数 第五版 北大数学系前线代小组 王萼芳 石生明 Z-Library.pdf
% IMG_2689中除以mu的条件订正为mu非零；IMG_2695第一列全零时取其余未知量全零。
% 极大无关组按非空组叙述；全零向量组的秩单独规定为0。
\RequirePackage{iftex}
\RequireLuaTeX
\documentclass[UTF8, fontset=fandol, a4paper, 12pt]{ctexart}

\usepackage[margin=2.5cm]{geometry}
\usepackage{amsmath, amssymb, amsthm}
\usepackage{enumitem}
\usepackage{fancyhdr}
\usepackage[hidelinks]{hyperref}

\newcommand{\va}{\boldsymbol\alpha}
\newcommand{\vb}{\boldsymbol\beta}
\newcommand{\vg}{\boldsymbol\gamma}
\newcommand{\vo}{\boldsymbol0}
\theoremstyle{definition}
\newtheorem{definition}{定义}[section]
\newtheorem{proposition}[definition]{命题}
\newtheorem{theorem}[definition]{定理}
\newtheorem{example}[definition]{例题}
\renewcommand{\proofname}{证明}
\setlist[enumerate]{leftmargin=2em, itemsep=0.2em, topsep=0.3em}
\setlength{\headheight}{16pt}
\pagestyle{fancy}
\fancyhf{}
\fancyhead[L]{\small 高等代数（一）2.1}
\fancyhead[R]{\small 线性相关性与向量组的秩}
\fancyfoot[C]{\thepage}
\renewcommand{\headrulewidth}{0.3pt}
\raggedbottom

\title{线性相关性与向量组的秩}
\author{Yuchen Zhang}
\date{整理日期：2026年11月10日}
\makeatletter
\renewcommand{\maketitle}{%
  \begin{center}
    {\LARGE\bfseries \@title\par}\vspace{0.6em}
    \@author\par\vspace{0.2em}
    {\small \@date\par}
  \end{center}
}
\makeatother

\begin{document}
\maketitle

\section{线性相关与线性无关}

以下 $P$ 为数域，所讨论的向量均属于 $P^n$，其中 $n\ge1$。
若存在 $c_1,\ldots,c_s\in P$，使
\[
 \vb=c_1\va_1+\cdots+c_s\va_s,
\]
则称 $\vb$ 可由向量组 $\va_1,\ldots,\va_s$ 线性表出。
特别地，取所有系数为零，总能得到零向量。线性相关性研究的是：
能否用一组不全为零的系数得到零向量。

\begin{definition}
设 $s\ge1$。若存在不全为零的 $c_1,\ldots,c_s\in P$，使
\[
 c_1\va_1+\cdots+c_s\va_s=\vo,
\]
则称向量组 $\va_1,\ldots,\va_s$ \textbf{线性相关}；否则称它\textbf{线性无关}。
因此，$\va_1,\ldots,\va_s$ 线性无关，当且仅当
\[
 c_1\va_1+\cdots+c_s\va_s=\vo
 \quad\Longrightarrow\quad c_1=\cdots=c_s=0.
\]
\end{definition}

\begin{proposition}
含有零向量的向量组线性相关；单个向量 $\va$ 组成的向量组线性无关，当且仅当 $\va\ne\vo$。
\end{proposition}
\begin{proof}
若组中第 $i$ 个向量为零，取 $c_i=1$，其余系数为零，就得到系数不全为零的零线性组合，
故该组线性相关。

若 $\va\ne\vo$，则它至少有一个分量 $a_j\ne0$。由 $c\va=\vo$ 得 $ca_j=0$，
从而 $c=0$，故单个向量 $\va$ 线性无关。若 $\va=\vo$，由前一结论知它线性相关。
\end{proof}

\begin{proposition}
两个向量 $\va,\vb$ 线性相关，当且仅当其中一个是另一个的数倍。
\end{proposition}
\begin{proof}
若 $\va,\vb$ 线性相关，则存在不全为零的 $\lambda,\mu\in P$，使
$\lambda\va+\mu\vb=\vo$。若 $\lambda\ne0$，则
$\va=-(\mu/\lambda)\vb$；若 $\lambda=0$，则 $\mu\ne0$，从而
$\vb=-(\lambda/\mu)\va$。

反之，若 $\va=c\vb$，则 $\va-c\vb=\vo$，其中 $\va$ 的系数为 $1$，故两向量相关。
若 $\vb=c\va$，同理可证。
\end{proof}
当一个向量为零、另一个非零时，只有零向量能由另一个向量线性表出，
因此命题中的“其中一个”不能固定为某一个。

\begin{proposition}
线性无关向量组的任意非空子组仍线性无关；一个向量组若含有线性相关的子组，则它线性相关。
\end{proposition}
\begin{proof}
若子组存在系数不全为零的零线性组合，给其余向量补上系数 $0$，
就得到整个向量组的一个系数不全为零的零线性组合。因此，相关子组所在的大组也相关。
取逆否命题，即得无关组的任意非空子组无关。
\end{proof}

\begin{theorem}\label{th:representation}
设 $s\ge2$。向量组 $\va_1,\ldots,\va_s$ 线性相关，当且仅当其中某个向量
可由其余向量线性表出。
\end{theorem}
\begin{proof}
若 $\sum_{j=1}^s c_j\va_j=\vo$，且系数不全为零，选取 $c_i\ne0$，则
\[
 \va_i=-\sum_{j\ne i}\frac{c_j}{c_i}\va_j.
\]
因此 $\va_i$ 可由其余向量线性表出。

反之，若 $\va_i=\sum_{j\ne i}d_j\va_j$，则
\[
 -\va_i+\sum_{j\ne i}d_j\va_j=\vo.
\]
其中 $\va_i$ 的系数为 $-1$，故原向量组线性相关。
\end{proof}

\section{线性相关性的例题}

\begin{example}
证明 $P^n$ 中的标准单位向量
\[
 \boldsymbol e_1=(1,0,\ldots,0)^{\mathsf T},\quad
 \boldsymbol e_2=(0,1,\ldots,0)^{\mathsf T},\quad\ldots,\quad
 \boldsymbol e_n=(0,0,\ldots,1)^{\mathsf T}
\]
线性无关。
\end{example}
\begin{proof}[解]
设 $c_1\boldsymbol e_1+\cdots+c_n\boldsymbol e_n=\vo$，则
\[
 (c_1,\ldots,c_n)^{\mathsf T}=\vo,
\]
故 $c_1=\cdots=c_n=0$。由定义，标准单位向量组线性无关。
\end{proof}

\begin{example}\label{ex:three}
判断下列向量组的线性相关性：
\[
 \va_1=(2,-1,4,-1)^{\mathsf T},\quad
 \va_2=(2,-1,3,1)^{\mathsf T},\quad
 \va_3=(4,-2,5,4)^{\mathsf T}.
\]
\end{example}
\begin{proof}[解]
令 $c_1\va_1+c_2\va_2+c_3\va_3=\vo$，按分量展开得
\[
 \begin{cases}
 2c_1+2c_2+4c_3=0,\\
 -c_1-c_2-2c_3=0,\\
 4c_1+3c_2+5c_3=0,\\
 -c_1+c_2+4c_3=0.
 \end{cases}
\]
第四式减去第二式得 $2c_2+6c_3=0$，故 $c_2=-3c_3$。
代入第二式得 $c_1=c_3$。取 $c_3=1$，得到 $(c_1,c_2,c_3)=(1,-3,1)$，
它也满足第一、三式。因此
\[
 \va_1-3\va_2+\va_3=\vo,
\]
而系数不全为零，故三个向量线性相关。
\end{proof}

\begin{example}
设 $\va_1,\va_2,\va_3$ 线性无关。证明：
\begin{enumerate}
 \item $\va_1+\va_2,\ \va_2+\va_3,\ \va_3+\va_1$ 线性无关；
 \item $\va_1-\va_2,\ \va_2-\va_3,\ \va_3-\va_1$ 线性相关。
\end{enumerate}
\end{example}
\begin{proof}[解]
对于第一组，设
\[
 c_1(\va_1+\va_2)+c_2(\va_2+\va_3)+c_3(\va_3+\va_1)=\vo.
\]
展开并按 $\va_1,\va_2,\va_3$ 合并，得到
\[
 (c_1+c_3)\va_1+(c_1+c_2)\va_2+(c_2+c_3)\va_3=\vo.
\]
由原向量组无关，得
\[
 c_1+c_3=0,\qquad c_1+c_2=0,\qquad c_2+c_3=0.
\]
前两式给出 $c_2=c_3=-c_1$，代入第三式得 $-2c_1=0$，故 $c_1=c_2=c_3=0$。
因此第一组线性无关。

对于第二组，有
\[
 (\va_1-\va_2)+(\va_2-\va_3)+(\va_3-\va_1)=\vo.
\]
系数 $(1,1,1)$ 不全为零，故第二组线性相关。
\end{proof}

\section{等价向量组}

\begin{definition}
设 $\mathcal A=(\va_1,\ldots,\va_s)$、$\mathcal B=(\vb_1,\ldots,\vb_t)$
是两个向量组。若每个 $\va_i$ 均可由 $\mathcal B$ 线性表出，则称
$\mathcal A$ 可由 $\mathcal B$ 线性表出。
若 $\mathcal A$ 与 $\mathcal B$ 可以互相线性表出，则称它们\textbf{等价}，记为
$\mathcal A\sim\mathcal B$。
\end{definition}
等价向量组所含向量个数可以不同。例如
$(\va_1,\va_2)$ 与 $(\va_1,\va_2,\va_1+\va_2)$ 等价。

\begin{proposition}
向量组的等价关系具有自反性、对称性和传递性。
\end{proposition}
\begin{proof}
\textbf{自反性：}每个向量都可由它所在的向量组线性表出，故 $\mathcal A\sim\mathcal A$。

\textbf{对称性：}若 $\mathcal A\sim\mathcal B$，则双方可以互相线性表出，
故 $\mathcal B\sim\mathcal A$。

\textbf{传递性：}设 $\mathcal A\sim\mathcal B$、$\mathcal B\sim\mathcal C$，
其中 $\mathcal C=(\vg_1,\ldots,\vg_u)$。由线性表示的定义，可写成
\[
 \va_i=\sum_{j=1}^{t}c_{ij}\vb_j,\qquad
 \vb_j=\sum_{k=1}^{u}d_{jk}\vg_k.
\]
代入并交换有限求和次序，得到
\[
 \va_i=\sum_{k=1}^{u}\left(\sum_{j=1}^{t}c_{ij}d_{jk}\right)\vg_k.
\]
因此 $\mathcal A$ 可由 $\mathcal C$ 表出。同理，由 $\mathcal C$ 可由 $\mathcal B$ 表出、
$\mathcal B$ 可由 $\mathcal A$ 表出，得 $\mathcal C$ 可由 $\mathcal A$ 表出。
故 $\mathcal A\sim\mathcal C$。
\end{proof}

\section{齐次方程组与向量个数的比较}

\begin{theorem}\label{th:homogeneous}
数域 $P$ 上的齐次线性方程组若方程个数 $m$ 小于未知量个数 $n$，则必有非零解。
\end{theorem}
\begin{proof}
对方程个数 $m$ 作归纳，证明结论对每个 $n>m$ 成立。
当 $m=0$ 时，未知量没有约束，取 $x_1=1$，其余未知量为零，即得非零解。

设结论对 $m-1$ 个方程成立，考虑 $m$ 个方程、$n>m$ 个未知量的齐次方程组
\[
 \sum_{j=1}^{n}a_{ij}x_j=0\qquad(i=1,\ldots,m).
\]
若第一列系数全为零，取 $x_1=1$、$x_2=\cdots=x_n=0$，即得非零解。

若第一列有非零系数，交换方程次序可设 $a_{11}\ne0$。
将第一式除以 $a_{11}$，再用第一式消去其余各式中的 $x_1$，得到同解方程组
\[
 \begin{cases}
 x_1+a'_{12}x_2+\cdots+a'_{1n}x_n=0,\\
 a'_{22}x_2+\cdots+a'_{2n}x_n=0,\\
 \hspace{5em}\vdots\\
 a'_{m2}x_2+\cdots+a'_{mn}x_n=0.
 \end{cases}
\]
后 $m-1$ 个方程有 $n-1$ 个未知量，且 $m-1<n-1$。
由归纳假设，存在不全为零的 $x_2,\ldots,x_n$ 满足这些方程。
再取
\[
 x_1=-a'_{12}x_2-\cdots-a'_{1n}x_n,
\]
即满足第一式。所得 $n$ 个分量不全为零，故为原方程组的非零解。归纳完成。
\end{proof}

\begin{theorem}\label{th:count}
若向量组 $\va_1,\ldots,\va_s$ 可由向量组 $\vb_1,\ldots,\vb_t$ 线性表出，且 $s>t$，
则 $\va_1,\ldots,\va_s$ 线性相关。
\end{theorem}
\begin{proof}
设
\[
 \va_i=\sum_{j=1}^{t}a_{ji}\vb_j\qquad(i=1,\ldots,s).
\]
考虑关于 $c_1,\ldots,c_s$ 的齐次方程组
\[
 \sum_{i=1}^{s}a_{ji}c_i=0\qquad(j=1,\ldots,t).
\]
因为方程个数 $t$ 小于未知量个数 $s$，由定理\ref{th:homogeneous}，
它有不全为零的解 $(c_1,\ldots,c_s)$。于是
\[
 \sum_{i=1}^{s}c_i\va_i
 =\sum_{i=1}^{s}c_i\sum_{j=1}^{t}a_{ji}\vb_j
 =\sum_{j=1}^{t}\left(\sum_{i=1}^{s}a_{ji}c_i\right)\vb_j=\vo.
\]
故 $\va_1,\ldots,\va_s$ 线性相关。
\end{proof}

\begin{proposition}\label{prop:independent-count}
若 $s$ 个线性无关向量可由 $t$ 个向量线性表出，则 $s\le t$。
两个等价的线性无关向量组含有相同个数的向量。
\end{proposition}
\begin{proof}
第一个结论是定理\ref{th:count}的逆否命题。
若两个等价的线性无关向量组分别含 $s$ 个和 $t$ 个向量，
由双方可以互相表出，分别得到 $s\le t$ 和 $t\le s$，故 $s=t$。
\end{proof}

\begin{proposition}
$P^n$ 中任意多于 $n$ 个向量组成的向量组线性相关。
\end{proposition}
\begin{proof}
任意 $\va=(a_1,\ldots,a_n)^{\mathsf T}\in P^n$ 均可写成
$\va=\sum_{j=1}^{n}a_j\boldsymbol e_j$。
因此 $P^n$ 中的任意向量组都可由 $n$ 个标准单位向量表出。
由定理\ref{th:count}，多于 $n$ 个向量的组必线性相关。
\end{proof}

\section{极大线性无关组}

\begin{definition}
设向量组 $\mathcal A$ 至少含有一个非零向量。它的一个非空子组 $\mathcal B$
称为 $\mathcal A$ 的\textbf{极大线性无关组}，若：
\begin{enumerate}
 \item $\mathcal B$ 线性无关；
 \item 向 $\mathcal B$ 添加 $\mathcal A$ 中任何一个未选向量，所得组都线性相关。
\end{enumerate}
\end{definition}

\begin{theorem}[极大无关组的等价刻画]\label{th:maximal}
原向量组的一个非空子组是极大线性无关组，当且仅当它线性无关，且原组的每个向量都可由它线性表出。
\end{theorem}
\begin{proof}
设 $\mathcal B=(\vb_1,\ldots,\vb_r)$ 为极大无关组。
任取原组中未选的向量 $\va$，将它加入 $\mathcal B$ 后所得向量组线性相关。
故存在不全为零的系数，使
\[
 c_1\vb_1+\cdots+c_r\vb_r+d\va=\vo.
\]
若 $d=0$，由 $\mathcal B$ 无关得 $c_1=\cdots=c_r=0$，矛盾。
故 $d\ne0$，从而
\[
 \va=-\frac{c_1}{d}\vb_1-\cdots-\frac{c_r}{d}\vb_r.
\]
$\mathcal B$ 中的向量当然也可由自身表出，所以原组的每个向量均可由 $\mathcal B$ 表出。

反之，设 $\mathcal B$ 无关且能表出原组的每个向量。添加任一未选向量 $\va$ 后，
由 $\va$ 可由 $\mathcal B$ 表出及定理\ref{th:representation}，所得组相关。
因此 $\mathcal B$ 是极大无关组。
\end{proof}

\begin{theorem}[存在性]
有限向量组只要含有非零向量，就存在极大线性无关组。
\end{theorem}
\begin{proof}
先选取一个非零向量，它单独组成线性无关组。
若原组每个向量均可由已选组表出，则由定理\ref{th:maximal}，已选组就是极大无关组。
否则选取一个不能由已选组表出的向量并加入。

加入后的组仍无关：若它相关，由已选组无关，零线性组合中新加入向量的系数必非零，
从而该向量可由已选组表出，与选取条件矛盾。
重复上述过程，每一步增加一个原组中的向量。由于原组有限，过程必在有限步后停止，
停止时所得组无关且能表出原组的所有向量，因此是极大无关组。
\end{proof}

\begin{theorem}
一个向量组的任意两个极大线性无关组都等价，且含有相同个数的向量。
\end{theorem}
\begin{proof}
每个极大无关组都是原组的子组，又能表出原组中的每个向量，因此它与原组等价。
由等价关系的传递性，任意两个极大无关组等价。
它们又都线性无关，由命题\ref{prop:independent-count}，所含向量个数相同。
\end{proof}

\begin{example}
求例\ref{ex:three}中向量组的极大线性无关组。
\end{example}
\begin{proof}[解]
由 $\va_1-3\va_2+\va_3=\vo$，三个向量线性相关。
其中任意两个都不成比例：$\va_1,\va_2$ 的第一分量相同而第三分量不同；
$\va_3$ 的第一分量是前两者的 $2$ 倍，但其第三分量既不是 $4$ 的 $2$ 倍，也不是 $3$ 的 $2$ 倍。
因此任意两个都线性无关。

选取任意两个向量后，添加余下的向量都得到原来的相关组，故
\[
 (\va_1,\va_2),\qquad(\va_1,\va_3),\qquad(\va_2,\va_3)
\]
均为极大线性无关组。
\end{proof}

\section{向量组的秩}

\begin{definition}
含有非零向量的向量组，其极大线性无关组所含向量的个数称为这个向量组的\textbf{秩}。
全由零向量组成的向量组，其秩规定为 $0$。
\end{definition}
极大无关组可以不唯一，但它们含有相同个数的向量，故秩的定义与极大无关组的选择无关。
例如例\ref{ex:three}中的向量组秩为 $2$。

\begin{proposition}
设 $\mathcal A$ 是含 $s\ge1$ 个向量的向量组，则
\[
 \mathcal A\text{ 线性无关}\ \Longleftrightarrow\ \operatorname{rank}\mathcal A=s,
 \qquad
 \mathcal A\text{ 线性相关}\ \Longleftrightarrow\ \operatorname{rank}\mathcal A<s.
\]
\end{proposition}
\begin{proof}
若 $\mathcal A$ 无关，则它本身就是极大无关组，故秩为 $s$。
若秩为 $s$，则一个极大无关组含有原组全部 $s$ 个向量，故原组无关。
由于极大无关组是原组的子组，其向量个数不超过 $s$，全零组的秩也小于 $s$。
因此第二个等价关系由第一个等价关系取否定得到。
\end{proof}

\begin{proposition}
等价向量组的秩相同。
\end{proposition}
\begin{proof}
设 $\mathcal A\sim\mathcal B$。若 $\mathcal A$ 全由零向量组成，则 $\mathcal B$ 中的每个向量
可由零向量表出，故也全为零，双方秩均为 $0$。

若 $\mathcal A$ 含有非零向量，则 $\mathcal B$ 也含有非零向量。
分别取双方的极大无关组 $\mathcal A'$、$\mathcal B'$，有
\[
 \mathcal A'\sim\mathcal A\sim\mathcal B\sim\mathcal B'.
\]
因此 $\mathcal A'$ 与 $\mathcal B'$ 是等价的线性无关向量组，所含向量个数相同，
故 $\mathcal A$ 与 $\mathcal B$ 的秩相同。
\end{proof}

\begin{example}
说明两个秩相同的向量组未必等价。
\end{example}
\begin{proof}[解]
在 $P^2$ 中，分别取单向量组
\[
 \mathcal A=\bigl((1,0)^{\mathsf T}\bigr),\qquad
 \mathcal B=\bigl((0,1)^{\mathsf T}\bigr).
\]
两个向量均非零，所以两组的秩都是 $1$。
但 $(0,1)^{\mathsf T}$ 的任意数倍第一分量都为零，不可能等于 $(1,0)^{\mathsf T}$，
故两组不等价。
\end{proof}

\section{课后作业}

\subsection{习题 III 第13题：由单位向量的表示证明线性无关}

设 $\va_1,\ldots,\va_n$ 是 $P^n$ 中的一组向量，已知标准单位向量
$\boldsymbol e_1,\ldots,\boldsymbol e_n$ 都可以由它们线性表出。
证明：$\va_1,\ldots,\va_n$ 线性无关。

\begin{proof}
一方面，由题设，$\boldsymbol e_1,\ldots,\boldsymbol e_n$ 可由 $\va_1,\ldots,\va_n$ 线性表出。
另一方面，$P^n$ 中的任意向量都可以由标准单位向量组线性表出，因而每个 $\va_i$ 也可以。
所以这两个向量组等价，具有相同的秩。

标准单位向量组线性无关，所含向量个数为 $n$，故它的秩为 $n$。
因此 $\va_1,\ldots,\va_n$ 的秩也是 $n$；它恰有 $n$ 个向量，由线性无关与秩的关系，
这组向量线性无关。
\end{proof}
本题题号按北大《高等代数》第六版标注。

\end{document}
